\answersfor{ch:shipping}

\begin{solution}{ex:shipping:1}
Compare first characters: \texttt{a} is 97, \texttt{h} is 104. Since $97 < 104$, \texttt{a\_spectral\_stem} sorts first and
\texttt{hyperspectral\_segformer\_mae} last. The dictionary keeps the last value written for ``indradhanu'', so it routes to
v0.2.0. Directory names deciding routing is exactly why the codenames must differ: the same two models would route
differently depending only on what a folder happened to be called.
\end{solution}

\begin{solution}{ex:shipping:2}
\begin{itemize}
\item 1400\,nm: $\abs{1400 - 1385} = 15$, $\abs{1460 - 1400} = 60$; minimum 15, and $15 \le 15$: kept.
\item 1420\,nm: $\abs{1420 - 1385} = 35$, $\abs{1460 - 1420} = 40$; minimum 35 $> 15$: dropped.
\item 1445\,nm: $\abs{1445 - 1385} = 60$, $\abs{1460 - 1445} = 15$; minimum 15: kept.
\end{itemize}
The rule keeps a 15\,nm fringe on each side of every recorded wavelength; a gap wider than 30\,nm in the record leaves
bands in its middle unread.
\end{solution}

\begin{solution}{ex:shipping:3}
Every band must pass \eqref{eq:shipping:allclose}; the tightest tolerance is at the smallest wavelength.
\begin{align*}
\text{at 460 nm:}\quad & 10^{-8} + 10^{-5} \times 460 = 0.0046\,\text{nm}\\
\text{at 2450 nm:}\quad & 10^{-8} + 10^{-5} \times 2450 = 0.0245\,\text{nm}
\end{align*}
A shift of 0.003\,nm is below 0.0046 at 460\,nm and below every larger tolerance, so \texttt{np.allclose} is true: the list
counts as the grid (and a native-band model would be refused). A shift of 0.03\,nm exceeds 0.0046 at 460\,nm (and even
0.0245 at 2450\,nm), so \texttt{np.allclose} is false: the list counts as native.
\end{solution}

\begin{solution}{ex:shipping:4}
\begin{align*}
100 \times 4 / 26 &= 400 / 26 = 15.3846\ldots\\
\text{rounded to 3 places} &= 15.385
\end{align*}
For $n = 200$: $n - T = 200 - 128 = 72$. The next multiple of 64 at or above 72 is 128, so the pad is $128 - 72 = 56$ (the
code's $(-72) \bmod 64 = 56$). The padded side is $200 + 56 = 256$, and windows start at $0, 64, 128$ (the last start is
$256 - 128 = 128$).
\end{solution}

\begin{solution}{ex:shipping:5}
Seven residuals: four land at $\pm 0.02$ and three cloud at $\pm 0.06$ magnitude.
\begin{align*}
\text{squares:}\quad & 4 \times 0.0004 + 3 \times 0.0036 = 0.0016 + 0.0108 = 0.0124\\
\text{variance:}\quad & 0.0124 / 7 = 0.001771\\
\text{spread:}\quad & \sqrt{0.001771} = 0.0421\\
\text{target:}\quad & 0.08 / 0.0421 = 1.90 \text{ spreads}
\end{align*}
More cloud in the fit widens the spread and lowers the target's score (from 2.09 to 1.90). Which pixels enter the fit
changes every score, which is why the product must use the order the benchmark measured.
\end{solution}

\begin{solution}{ex:shipping:6}
The action fails at run time. The default band filter resamples to the 165-band grid, so V2's grid check raises ``reads
native bands: run band\_filter\_apply with use\_common\_wavelength\_grid=false''. Instead, run two scoring actions: v0.2.0
on a default band filter, and V2 on a second band filter with \texttt{use\_common\_wavelength\_grid=false}.
\end{solution}
